% This file should be rename with the speaker's last name. (e.g : smith)

% Write the author names in the order you wish them to appear
% and underline the one who will give the talk.

% Only the speaker's email address is required.

\begin{abstract_online}{Symbolic calculation behind floating-point arithmetic using CAS}{%
        \underline{W{\l}odzimierz Wojas}$^{1}$, \underline{Jan Krupa}$^{1}$}{[jan\_krupa@sggw.pl]
        }{%
        $^1$ Department of Applied Mathematics, Warsaw University of Life Sciences (SGGW), Warsaw, Poland}




In this talk we would like to present using CAS, some examples of symbolic calculations which lie behind calculations in floating-point arithmetic (with double precision). Each operation in floating-point arithmetic is performed according to a precise-symbolic algorithm. In spite of the fact that floating point arithmetic is based on symbolic operations, it gives approximate results with the exception of adding, subtracting, multiplying and dividing powers of two, adding, subtracting, and multiplying sums of powers of two, adding, subtracting and multiplying of integers.
We will present in this talk the simple examples in Mathematica and wxMaxima that result of operations depends on the interpretation of the used input data (numbers) by CAS functions(such as Solve, Limit, Det, solve, limit, det) – as symbolic or approximate and on used by these CAS functions algorithms.


\textbf{Keywords} \newline{} Higher education,Floating Point Arithmetic, Application of CAS, Mathematica, Mathematical didactics




\textbf{References}\newline{}
[1] \textsc{Timothy Sauer},  \emph{Numerical Analysis.} 3rd Edition, Pearson,  2017 \newline{}
[2] IEEE 754 Floating Point Standard: https://en.wikipedia.org/wiki/IEEE\_754 \newline{}
[3] \textsc{D. Goldberg},  “What Every Computer Scientist Should Know about Floating Point
Arithmetic.’’ ACM Computing Surveys 23, 5–48, 1991 \newline{}
[4] \textsc{W. Stallings}, \emph{Computer Organization and Architecture}, 6th ed. Prentice Hall, Upper Saddle River, NJ, 2003.
\end{abstract_online}